\documentclass[11pt]{article}

% arXiv cs.MS (cross-list cs.LO) submission.
\usepackage[margin=1.1in]{geometry}
\usepackage[T1]{fontenc}
\usepackage[utf8]{inputenc}
\usepackage{amsmath,amssymb,amsthm}
\usepackage{booktabs}
\usepackage{microtype}
\usepackage[hidelinks]{hyperref}
\usepackage{url}
\usepackage{listings}
\usepackage{xcolor}

\lstdefinelanguage{lean}{
  keywords={structure, inductive, theorem, def, lemma, instance, namespace,
            end, import, where, by, fun, match, with, deriving},
  sensitive=true,
  comment=[l]{--},
  morestring=[b]",
}
\lstset{
  language=lean,
  basicstyle=\small\ttfamily,
  keywordstyle=\bfseries,
  commentstyle=\itshape\color{gray},
  columns=fullflexible,
  breaklines=true,
  frame=single,
  framerule=0.3pt,
  xleftmargin=1em,
}

\newtheorem{theorem}{Theorem}
\newtheorem{definition}{Definition}

% Release-mode guard: release builds hard-fail on missing rendered
% fragments; draft builds show placeholders. There are no generated files —
% only artifacts with receipts: authority comes from the .mfact/artifacts.toml
% ledger, not from paths or headers. Volatile standing claims live only in
% ledgered, ggen-rendered fragments.
\newif\ifreleasebuild
\releasebuildtrue
% GENERATED FILE — DO NOT EDIT BY HAND
% Source: release-manifest.json + standing.env via the quadrature graph
% Generator: ggen (quadrature-pack) — ggen renders; Lean admits; mfact certifies.
% Standing: GENERATED_FROM_CERTIFIED_MANIFEST
\newcommand{\ReleaseTag}{v26.7.7-procint-certified}
\newcommand{\ReleaseHash}{942facf32d48cd1a26c0f06b9396c6c150ab4d95d601bd090a8e1b9e7ef2d434}
\newcommand{\ReleaseRun}{945bfca}
\newcommand{\SorryCount}{0}
\newcommand{\AdmitCount}{0}
\newcommand{\KernelTheorems}{197}
\newcommand{\TotalDecls}{397}
\newcommand{\StatedCount}{7}
\newcommand{\AxiomAuditStatus}{PASS}
\newcommand{\FixtureStatus}{PASS}
\newcommand{\CertifiedStatus}{PASS}
\newcommand{\QuadratureStatus}{PASS}
\newcommand{\CrownJewelStatus}{PROVEN}


\title{MathProofOps:\\Receipted Mathematical Manufacturing for Process Intelligence\\
\large From Public Ontology to Lean-Checked Standing}
\author{Sean Chatman}
\date{}

\begin{document}
\maketitle

\begin{abstract}
We present \textsf{MathProofOps}, a discipline for lawful mathematical
artifact manufacture, and \textsf{mfact}, its Lean~4 manufacturing
framework. The governing law is \(A = \mu_B(O^*_B)\) and \(R_B =
\mathrm{receipt}(A)\): inside a declared boundary \(B\), raw observations
and generated candidates have no standing. Only admitted observations
\(O^*_B\) may participate in manufacture; only replayable receipts \(R_B\)
prove consequence. This separates contact from consequence. Humans,
templates, and LLM agents may produce candidates; only the Lean kernel,
Lake build, axiom audit, sorry census, negative fixtures, and release
manifest can confer standing.

As the first manufactured vertical we introduce \textsf{procint}, a Lean~4
library for process intelligence: Petri nets with finitely supported
markings, workflow nets with soundness split into independent clauses,
event logs and alignment-based conformance, and object-centric event logs
following OCEL~2.0. The library is manufactured from a declared obligation
graph: an RDF/Turtle catalog of structures, theorem surfaces, module
dependencies, fixtures, citations, and registry surfaces is projected
through SPARQL-in-Tera templates by \textsf{ggen}, then admitted or
refused by Lean/Lake. Every failed admission is a typed refusal rather
than a silent downgrade; every admitted declaration is recorded with its
hash, axiom footprint, fixture status, and proven/stated standing. The
contribution is not merely that \textsf{procint} exists; it is that
\textsf{procint} is the first manufactured vertical of MathProofOps, and
\textsf{mfact} is the Lean/Lake manufacturing apparatus that makes this
vertical receiptable.

A second rail extends the manufacturing graph from domain semantics to
implementation correspondence. Using Aeneas/Charon, bounded Rust cores from
\textsf{wasm4pm}-style process execution can be extracted into
proof-facing semantic images. These images do not prove Rust correct by
extraction alone; rather, they create correspondence obligations between
implementation objects and admitted \textsf{procint} semantics under
declared abstraction functions. This gives process intelligence a proof
class it otherwise lacks: kernel-checked implementation-correspondent
process law, not merely ``model proved, implementation tested.''

The result is not a claim that arbitrary generated mathematics is
automatically correct. Rice-style undecidability is respected rather than
ignored: MathProofOps replaces the impossible question ``is this arbitrary
generated artifact semantically correct?'' with the bounded, replayable
question ``did this candidate pass the declared admission gate inside
boundary \(B\)?'' More generally, public ontologies supply a domain's
nouns; MathProofOps supplies standing. The method generalizes from process
intelligence to any ontology-bearing field whose artifacts can be
projected, admitted, receipted, and replayed. The evaluation section of
this paper is generated from the release manifest rather than written by
hand, and a kernel-checked \emph{Standing Quadrature} witness closes the
cross-product between the declaration catalog, the admitted corpus, the
release manifest, the process evidence, and this paper's claims. We report
exactly what was proven, what remains stated, what was refused, and what
the release boundary excludes.
\end{abstract}

\section{Introduction}

Formalization today often resembles hand-fitted manufacturing: every
theorem, lemma, import, and repair is shaped locally by expert craft. The
result can be excellent, but it is not yet industrial. Mathematical
manufacturing asks a different question: what must be standardized so that
formal domains can be produced as interchangeable parts? In
\textsf{mfact}, the interchangeable part is not merely a Lean declaration.
It is a declaration plus provenance, admission status, axiom footprint,
refusal mode, hash, manifest position, and replay path. A theorem without
this surrounding receipt is hand-fitted metal; an admitted declaration
with its release evidence is an interchangeable mathematical part. The
manufactured object is not syntax --- it is \emph{standing}.

Recent work uses LLMs to draft formal content, which raises a question
the kernel alone does not settle: what standing does generated material
have? Kernel acceptance establishes that a term inhabits its stated type;
it does not establish that the stated type is the intended one, that the
statement is not vacuously weakened, or that the artifact as a whole is a
faithful rendering of a target theory. Meek et
al.~\cite{meek2026formalizing} document these failure modes in detail:
compilation masks incomplete multi-part statements, added weakening
hypotheses, and parameter restrictions. Our position is compressed by one
rule: \emph{contact is not consequence}. LLM contact with a theorem is
not consequence. A model may propose a proof term, a definition, or a
module, but the proposal has no standing until admitted. \textsf{mfact}
formalizes this separation: contact produces a candidate; admission
produces consequence; the manifest records the receipt.

Our position is that the two concerns should be mechanically separated:

\begin{itemize}
\item \textbf{Generation is unprivileged.} Anything --- a template engine, a
  human, an LLM --- may produce a \emph{candidate}. Candidates carry
  provenance and a content hash but no standing.
\item \textbf{Standing is admission.} A candidate acquires standing only by
  passing a fail-closed gate: the Lean kernel accepts it, it contains no
  \texttt{sorry}, and its axiom footprint is pinned to an authorized set.
  Every failure is a typed refusal, enumerated in a closed inductive type.
\item \textbf{Claims are scoped.} A release manifest records, per
  declaration, whether it is \emph{proven} (kernel-admitted, axiom-audited)
  or merely \emph{stated} (the formal statement exists; the proof
  obligation is open). The paper's own evaluation is rendered from the
  manifest, not asserted in prose.
\end{itemize}

We call this discipline \emph{mathematical manufacturing}: the pipeline is
the product as much as the library is. The statement--proof split, the typed
refusals, and the manifest are what make an overnight, largely unattended
run reportable at all.

\paragraph{Scope.} This paper does not claim to formalize all of process
intelligence completely, nor does it claim the manufacturing pipeline is
applicable beyond the case study reported here. It introduces a
mathematical manufacturing pipeline and demonstrates it by generating a
first Lean~4 process-intelligence foundation from an existing typed
implementation inventory. Table~\ref{tab:stages} makes the pipeline
concrete: each stage names the artifact it produces, not a metaphor.

\begin{table}[h]
\centering
\small
\begin{tabular}{ll}
\toprule
Manufacturing stage & Concrete artifact \\
\midrule
Source inventory & typed Rust implementation (wasm4pm-compat) \\
Rendering & \textsf{ggen} SPARQL-in-Tera templates (render Lean modules from the TTL catalog) \\
Admission & Lean module creation (\texttt{lake env lean}) \\
Inspection & \texttt{lake build} \\
Refusal & typed \texttt{Refusal}, negative fixtures \\
Evidence & process trace, release manifest \\
Certification & \texttt{CertifiedRelease} object, axiom audit \\
Publication & generated paper evaluation section \\
\bottomrule
\end{tabular}
\caption{Mathematical manufacturing stages and their artifacts.}
\label{tab:stages}
\end{table}

This paper makes four bounded contributions.

First, it defines \emph{mathematical manufacturing} as a typed admission
discipline for formal artifacts. Generated output, human patches,
templates, and agent proposals are treated uniformly as candidates.
Standing is acquired only through admission: kernel acceptance, no-sorry
census, axiom audit, negative fixtures, manifest completeness, and
replayable receipt (Section~\ref{sec:law}).

Second, it introduces \textsf{mfact}, a Lean~4 framework that makes this
discipline explicit through candidate, refusal, admission, manifest,
certified-release, and valid-objection types (Section~\ref{sec:mfact}).
The framework is intentionally small: it does not attempt to replace the
Lean kernel, but to make the pre-kernel and post-kernel standing surface
explicit.

Third, it introduces \textsf{procint}, the first manufactured vertical of
MathProofOps, produced by \textsf{mfact}: a process-intelligence library
covering Petri nets, workflow nets, event logs, conformance, and
OCEL-style object-centric logs (Section~\ref{sec:procint}). The novelty is
not only the domain coverage --- to our knowledge, based on a bounded
search of the Lean, Coq, and Isabelle/AFP library indexes at the time of
the run, we found no prior public mechanization of workflow-net soundness
or OCEL in these systems, a negative search result rather than a proof of
absence --- it is the coupling of domain formalization with an explicit
manufacturing receipt. A prior hand-written mechanization, if found, would
be adjacent but not equivalent unless it provides the same admission,
refusal, manifest, axiom, fixture, and replay surfaces.

Fourth, it reports a manifest-generated evaluation
(Section~\ref{sec:eval}) closed by a kernel-checked Standing Quadrature
witness. The paper's evaluation numbers are not prose claims; they are
rendered from the release manifest, and a generated Lean module
(\texttt{ProcInt.Release.Quadrature}) is refused by the kernel if the
declaration catalog, the audited corpus, and the manifest ever diverge.
This closes the loop from declaration to admission to receipt to replay,
and makes the paper itself an artifact of the manufacturing pipeline it
describes.

\section{The Receipted Manufacturing Law}
\label{sec:law}

The governing law of the pipeline is
\[
A = \mu_B(O^*_B), \qquad R_B = \mathrm{receipt}(A)
\]
where \(B\) is the declared boundary (the Lean packages, the declaration
catalog, the manifest schema, the source inventory, the audited theorem
set); \(O\) is raw observation --- candidate output from a human, an LLM,
a template, a repository, or an ontology fragment; \(O^*_B\) is admitted
observation inside \(B\): typed, checked, bounded, replayable; \(\mu_B\) is
the lawful manufacturing transform (\textsf{ggen} rendering, Lean
admission, Lake build, audit), parameterized by the boundary it operates
inside; \(A\) is an artifact with standing (an admitted theorem, module,
manifest entry, or release); and \(R_B\) is the receipt (build log,
manifest hash, axiom audit, negative fixture, replay evidence) that proves
\(A\) is a genuine consequence of \(\mu_B\), not an assertion about it.
Raw generated output is \(O\), not \(O^*\). Manufacture produces the
artifact; the receipt proves consequence; replay proves the receipt is
non-ornamental. \emph{The pipeline exists to refuse the gap.}

\paragraph{The Four Evidence Kinds.} We distinguish four classes of evidence within the manufacturing graph: (1) Formal evidence (Lean theorems and axiom audits); (2) Operational evidence (praxis-graphlaw benchmark and execution validation via rslab); (3) Correspondence evidence (Aeneas-extracted image to procint semantic equivalence proofs); and (4) Publication evidence (generated paper fragments and manifests).

In this architecture, \textsf{ggen} is not a code generator in the usual
sense. It is the projection half of the manufacturing function \(\mu\):
it maps a declared obligation graph into candidate Lean surfaces. Lean
supplies the admission half. Projection without admission is not
standing; admission without projection does not scale.

\paragraph{Rice boundary and manufactured decidability.}
Rice's theorem bars nontrivial semantic decision procedures over
arbitrary programs. The manufacturing response is not to evade this
boundary, but to change the object of judgment. MathProofOps does not
decide arbitrary semantic truth of arbitrary generated code. It decides
whether a bounded candidate has been admitted by a specific kernel, under
a declared axiom policy, with a complete manifest receipt. The
undecidable question ``is this arbitrary generated artifact
mathematically correct?'' is replaced by the decidable question ``did
this candidate pass the declared admission gate inside boundary \(B\)?''
Undecidability is not defeated; it is fenced, and bounded standing is
manufactured inside the fence.

\paragraph{Status is algebra, not prose.}
The release distinguishes a closed vocabulary of standing values:
\textsc{Declared}, \textsc{Extracted}, \textsc{Stated}, \textsc{Proven},
\textsc{Refused}, \textsc{Blocked}, \textsc{BuildBroken}, \textsc{Unknown},
\textsc{Unsupported}, and \textsc{Alive}. No transition from
\textsc{Stated} to \textsc{Proven} is implicit or inferred from prose;
\textsc{Extracted} is not proof; \textsc{Unknown} is not admitted;
\textsc{Unsupported} is not the same as \textsc{Refused}. A generated
artifact remains a candidate, whatever its origin, until an explicit
admission step assigns it one of these values.

\section{From Public Ontology to Domain Standing}
\label{sec:ontology-generalization}

The method is not specific to process intelligence. Process intelligence
is the first vertical because its public objects --- events, traces,
objects, transitions, guards, effects, conformance obligations, and
object-centric logs --- already carry close mathematical structure. The
underlying pattern is field-independent. Public ontologies supply a
domain's nouns; they do not supply standing. \textsf{ggen} projects those
nouns into candidate artifacts; Lean/Lake admits or refuses them;
manifests record algebraic status; receipts and replay preserve
consequence across time.

For a field \(F\) with public object language \(\mathcal{O}_F\), raw
observations \(O_F\), admitted observations \(O^*_F\), lawful
manufacturing transform \(\mu_F\), artifact \(A_F\), and receipt \(R_F\),
the MathProofOps law specializes exactly as in Section~\ref{sec:law}:
\[
A_F = \mu_F(O^*_F), \qquad R_F = \mathrm{receipt}(A_F).
\]
Public ontology supplies the field's nouns; MathProofOps supplies
standing. Finance has FIBO/XBRL/LEI; healthcare has FHIR/SNOMED/LOINC;
supply chain has GS1/EPCIS/PROV-O; manufacturing has SOSA/QUDT/ISA-95;
software has SPDX/SBOM and extracted implementation images. In each case
the platform claim is not that the ontology is new, but that
ontology-grounded artifacts become admitted, receipted, replayable objects
with standing. \textsf{procint} is this paper's demonstration vertical,
not the boundary of the method.

\section{Architecture: The Standing-Manufacturing Graph}
\label{sec:architecture}

The standing-manufacturing graph is organized around a five-rail architecture. Table~\ref{tab:rails} details each rail's source, admission boundary, and receipt form.

\begin{table}[h]
\centering
\small
\begin{tabular}{llll}
\toprule
Rail & Source & Admission Boundary & Receipt Form \\
\midrule
Formal & Lean catalog, corpus & Lean kernel, axiom policy & Audit log, print axioms \\
Process-law & mfact core definitions & Conformance obligs. & covers proof, ValidObjection \\
Benchmark & praxis execution & rslab schemas & JSON/TOML result receipt \\
Correspondence & Rust crate source & Aeneas extraction & build\_verif.py check \\
Paper & Manifest status & ggen sync & Manifest foldHash, regen-check \\
\bottomrule
\end{tabular}
\caption{The Five-Rail standing-manufacturing architecture.}
\label{tab:rails}
\end{table}

\section{Related Work}

\paragraph{Formal libraries and their pipelines.}
Mathlib~\cite{mathlib2020} is the reference point for large-scale
collaborative formalization in Lean; its quality regime is human review plus
continuous integration. CSLib~\cite{cslib2026} extends the Lean ecosystem
toward computer science, including labeled transition systems, traces, and
bisimulation, which \textsf{procint} builds on for its semantics layer. Lean~4
itself is described by de Moura and Ullrich~\cite{moura2021lean4}. The
closest neighbor to \textsf{mfact} is the agent pipeline of Meek et
al.~\cite{meek2026formalizing}, which formalizes numerical analysis with
coding agents and audits quality beyond kernel acceptance. We differ in
mechanism rather than diagnosis: where their audit is a post-hoc quality
assessment over agent output, \textsf{mfact} is a \emph{typed pipeline} ---
candidates, refusals, and admissions are Lean data types, the statement
corpus lives in an RDF catalog curated by LLM-assisted lanes (each
candidate kernel-verified before being recorded, rather than accepted on
the strength of agent output alone), and the proven/stated split is
enforced by the release gate rather than estimated afterward. The two
approaches are complementary: their semantic audits address exactly the
residual risk our pipeline does not eliminate, namely that a curated
statement, though kernel-admitted, is itself a faithless or vacuous
rendering of the intended theorem.

\paragraph{Lean 4 automation infrastructure.} \textsf{mfact}'s trust
boundary sits above an already-scrutinized kernel: Lean4Lean provides an
independent, kernel-in-Lean typechecker whose verification effort has
already located and fixed a soundness bug in the reference
implementation~\cite{lean4lean2024}. The two dominant patterns for wiring
language-model or solver automation into Lean~4 are in-process integration,
as in Lean Copilot's LLM inference over an FFI boundary~\cite{leancopilot2024},
and external bridging to automated theorem provers, as in
Lean-auto~\cite{leanauto2025}. \textsf{mfact} adopts neither: candidates from
any producer, language model or otherwise, are untrusted data admitted only
through the gate of Section~\ref{sec:mfact}, so no automation tool sits on
the trusted path regardless of how it is wired in. Infrastructure for
premise and declaration search~\cite{leanexplore2025} and for large
synthetic theorem corpora~\cite{mathlib4dataset2025} addresses problems
\textsf{procint} does not currently have, since its statement corpus is
curated from the RDF catalog rather than mined or synthesized.

\paragraph{Process intelligence.}
The mathematical canon \textsf{procint} formalizes is due to a small number
of primary sources. Petri-net structure theory --- state equation,
invariants, boundedness, liveness --- follows Murata's
survey~\cite{murata1989}. Workflow nets and their soundness are van der
Aalst's~\cite{vanderaalst1997}; the decidability and classification
landscape is surveyed in~\cite{vanderaalst2011soundness}, and the process
mining setting in~\cite{vanderaalst2011pm}. Alignment-based conformance
checking follows Carmona et al.~\cite{carmona2018}. Object-centric event
data follows the OCEL~2.0 specification~\cite{ocel2} and object-centric
Petri nets~\cite{ocpn2020}; partially ordered workflow models follow
POWL~\cite{powl2023}. None of these, to the best of our search, had a prior
mechanization in a proof assistant.

\paragraph{Code-to-proof extraction.} The correspondence rail of
Section~\ref{sec:correspondence} extracts an external Rust crate into Lean
via Aeneas~\cite{ho2022aeneas}, whose Charon front end~\cite{protzenko2024charon}
separates syntax recovery from semantic translation. Aeneas targets whole
Rust programs with a full memory model; our pilot uses it far more
narrowly, on a single dependency-free arithmetic module, and treats its
output the same way as any other candidate in this paper's pipeline ---
extraction produces a Lean statement, never a proof, and standing still
requires kernel admission of a hand-authored correspondence theorem
against that extraction. This rail addresses a proof class ordinary
process intelligence lacks: without it, a Lean process model may be proven
and a Rust executor may be tested, but their relationship remains prose,
naming, convention, and examples. Extraction changes the object class, so
that implementation and semantics can be related by a kernel-checked
theorem rather than by documentation.

\section{The mfact Framework}
\label{sec:mfact}

\textsf{mfact} is deliberately small: seven modules and a CLI. Its job is to make the pipeline's epistemics first-class Lean objects. mfact is not the whole system: it is the Lean/Lake standing apparatus inside a larger standing-manufacturing graph.

\paragraph{Candidates and refusals.}
A \texttt{Candidate} is content plus provenance (producer identity, content
hash); producers are untrusted by construction. \texttt{Refusal} is a closed
inductive type --- \texttt{kernelRejected}, \texttt{sorryPresent},
\texttt{unauthorizedAxiom}, \texttt{missingEvidence},
\texttt{malformedModel}, \texttt{invalidFixture} --- so a failed admission is
always a value that can be logged, hashed, and re-queued, never an untyped
error string. There is no default case: an unrecognized failure mode is
itself a refusal to classify, not a silent pass.

\paragraph{Admission with a covers proof.}
An \texttt{Admission} bundles the evidence for one candidate: the kernel
acceptance record, the sorry census, the axiom audit output, and a
\texttt{covers} field --- a proof term witnessing that the evidence entries
cover the obligations the gate imposes, so an admission cannot be assembled
with a missing check. The gate function \texttt{admit} is total: it returns
either an \texttt{Admission} or a \texttt{Refusal}, with no third outcome.

\paragraph{Manifest, release, objections.}
A \texttt{Manifest} (JSON-serializable via \texttt{ToJson}) lists artifacts
with hashes, axiom sets, fixture results, and the proven/stated status of
each declaration. A \texttt{CertifiedRelease} pairs a manifest with its
audit obligations. \texttt{ValidObjection R} is a closed inductive family
whose constructors are the only admissible objection forms against a
release \texttt{R} (an unaudited theorem, a missing hash, an unauthorized
axiom, a failing fixture); each constructor demands evidence. The
per-release theorem
\begin{center}
\texttt{theorem no\_valid\_objection : IsEmpty (ValidObjection R)}
\end{center}
is discharged by refuting each constructor against the corresponding checked
field of \texttt{R}. For the release described in this paper the theorem is
proven and, per \texttt{\#print axioms}, depends on no axioms at all --- not
even the three standard Lean/Mathlib axioms. The theorem's scope is exactly
its constructors: it says the \emph{enumerated} objection forms are refuted,
not that the release is beyond all criticism.

\section{procint: Process Intelligence as Process Law}
\label{sec:procint}

In this paper, process intelligence does not mean dashboarding over event
logs. It means a mathematical substrate for deciding, replaying, and
explaining process standing: what happened, what was admitted, what was
refused, which object relations were involved, which conformance
obligations passed, and which receipts authorize the resulting claim.
\textsf{procint} depends on Mathlib and CSLib and is organized as follows
(module map abridged):

\begin{center}
\small
\begin{tabular}{ll}
\toprule
Namespace & Content \\
\midrule
\texttt{ProcInt.Foundations} & metrics as \(\{q : \mathbb{Q} \mathrel{/\!/} 0 \le q \land q \le 1\}\), identifiers, LTS \\
\texttt{ProcInt.Petri} & nets, firing, reachability, state equation, invariants, \\
 & boundedness, liveness, OCPN, stochastic nets, unfoldings \\
\texttt{ProcInt.Workflow} & WF-nets, short-circuit construction, soundness \\
\texttt{ProcInt.Logs} & events, traces (monotone timestamps), event logs, XES \\
\texttt{ProcInt.Models} & process trees, POWL, DFGs, causal nets, BPMN, Declare \\
\texttt{ProcInt.Conformance} & moves, alignments, cost, fitness, token replay \\
\texttt{ProcInt.Ocel} & OCEL~2.0, E2O/O2O relations, flattening, lifecycles \\
\texttt{ProcInt.Ocpq} & object-centric process queries \\
\texttt{ProcInt.Analytics} & temporal, causality, correlation, cubes, streaming \\
\texttt{ProcInt.Registry} & algorithm and breed registries (cross-product tables) \\
\bottomrule
\end{tabular}
\end{center}

\paragraph{Petri nets.}
A net is places, transitions, and weighted flow over finite types; a marking
is a finitely supported function \(M : P \to_{f} \mathbb{N}\), which unifies
the weighted and (via richer codomains) colored cases and gives the marking
algebra of Murata~\cite{murata1989} --- firing, the state equation
\(M' = M + A^{\mathsf T} u\), and place/transition invariants --- as
\(\mathsf{Finsupp}\) lemmas with truncated subtraction.

\paragraph{Workflow nets and soundness.}
A \texttt{WfNet} carries a Petri net as a field (composition, not
extension) together with the source-place, sink-place, and connectivity
conditions of van der Aalst's Definition~6--7~\cite{vanderaalst1997}.
Following the observation that the classical soundness clauses are logically
independent, soundness is \emph{not} a single conjunction-shaped definition
but three separately named predicates --- option to complete, proper
completion, and no dead transitions --- with \texttt{Sound} as their
explicit conjunction. This keeps partial results honest: a lemma about one
clause is stated about that clause.

The crown-jewel target of the run is van der Aalst's soundness
characterization (Lemma~8~\cite{vanderaalst1997}): soundness of a workflow net
iff liveness and boundedness of its short-circuited net. The finite transition
hypothesis is necessary: the repository includes a Lean-admitted
infinite-transition countermodel, \verb|WfNet.infinite_transition_countermodel_sound_not_bounded|,
showing that the unbounded statement fails under infinite transitions. The
repaired crown theorem therefore states the equivalence under \([Finite T]\).

\paragraph{Conformance.}
Alignments follow Carmona et al.~\cite{carmona2018}: moves are synchronous,
log-only, model-only, or silent; costs induce optimal alignments; fitness is
a \(\mathsf{Between01}\) value by construction, so boundedness of the metric
is a type-level fact rather than a theorem obligation.

\paragraph{OCEL 2.0.}
Object-centric event logs follow the OCEL~2.0
specification~\cite{ocel2} (Definitions~1--2): typed events and objects,
event-to-object and object-to-object relations with qualifiers, and
attribute maps; flattening onto a classical log and object lifecycles are
derived operations. Object-centric Petri nets follow~\cite{ocpn2020}.

\paragraph{Canon status.} Table~\ref{tab:canon} states, per area, whether
\textsf{procint} v1 formalizes it, formalizes a seed instance of it, or
leaves it to future work. The claim of this paper is a first combinatorial
process-intelligence foundation, not a complete formalization of process
intelligence; per-declaration status (proven/stated) within each formalized
area is in the generated evaluation (Section~\ref{sec:eval}).

\begin{table}[h]
\centering
\small
\begin{tabular}{ll}
\toprule
Area & v1 status \\
\midrule
Transition systems, traces, reachability & formalized (Foundations, Petri) \\
Petri nets (marking, firing, invariants) & formalized \\
Workflow nets and soundness & formalized; short-circuit theorem partial (Sec.~\ref{sec:eval}) \\
Alignment-based conformance & formalized (moves, cost, fitness) \\
OCEL~2.0 event/object relations & formalized \\
Object-centric Petri nets & seed instance \\
Process trees, POWL, DFGs, BPMN, Declare & seed instances \\
Object-centric process queries (OCPQ) & seed instance \\
Temporal/causal/streaming analytics & seed instances \\
\bottomrule
\end{tabular}
\caption{Process-intelligence canon coverage in \textsf{procint} v1.}
\label{tab:canon}
\end{table}

\section{The Manufacturing Run}
\label{sec:run}

The run that produced \textsf{procint} has five moving parts.

\paragraph{The TTL catalog as canonical source.}
The canonical source of \textsf{procint} is not a tree of \texttt{.lean}
files but an RDF/Turtle declaration catalog (\texttt{procint:} vocabulary):
one record per declaration, carrying its Lean body (\texttt{leanCode}),
its module, its proof status, its axiom-audit expectation, and citations
back to the primary literature. Declaration bodies enter the catalog by
curation --- LLM-assisted lanes draft them and each candidate is
kernel-tested before being recorded --- so the catalog holds verified
candidates, not free-form generation. The cross-product surfaces
(registries of algorithms and model breeds) are fully data-driven rows of
the same graph.

\paragraph{SPARQL-in-Tera rendering.}
The \textsf{ggen} engine renders Lean modules from templates whose
frontmatter carries SPARQL \texttt{SELECT} queries against the catalog and
whose bodies are Tera templates over the result bindings. One template
yields one \texttt{.lean} file per \texttt{procint:Module}; separate
templates render the root import index, the registry tables, and the
axiom-audit file --- for every audited theorem, a \texttt{\#guard\_msgs in
\#print axioms} block pinned to \texttt{[propext, Classical.choice,
Quot.sound]}, so any transitive \texttt{sorryAx} or unauthorized axiom
breaks the build. The rendered corpus is a \emph{candidate artifact}: it is
not trusted because it was rendered, and \textsf{ggen} certifies nothing.
It acquires standing only through kernel admission. In one line:
\textsf{ggen} renders, Lean admits, \textsf{mfact} certifies.

\paragraph{LLM lanes as untrusted producers.}
Module families (Petri core, workflow, logs and conformance, OCEL,
models, analytics, fixtures) run as independent lanes with no barriers
between them. Each lane's agent hand-drafts the actual structures,
definitions, and proofs for its module family, but writes nothing directly
into the library: every piece of Lean code it proposes is first
kernel-tested via \texttt{lake env lean} on a scratch file, and only a
declaration that passes this check is recorded --- as a \texttt{leanCode}
string literal, together with its proof status --- in the lane's catalog
fragment. The catalog records what the lane authored and what the kernel
confirmed; it authors nothing itself. Statements the lane cannot prove ship
as \texttt{sorry}-bearing stubs marked \emph{stated} in the fragment.

\paragraph{Admission and the repair loop.}
An integrator loop repeatedly merges fragments, re-renders, and runs
\texttt{lake build}. Build errors are parsed and returned to the owning
lane as repair context, with bounded retries per stub and per lane;
exhausted stubs are downgraded --- the statement stays in the library, the
claim drops to \emph{stated}, and the declaration leaves the audited set.
Nothing is deleted and nothing passes silently: the terminal state of every
stub is either kernel-admitted or an explicit downgrade receipt.

\paragraph{Certification.}
The release gate rebuilds both packages, builds the axiom-audit files,
runs the negative fixtures (\texttt{\#guard\_msgs (error)} blocks, the Lean
analog of compile-fail tests), and emits the release manifest with content
hashes. A negative control --- injecting a \texttt{sorry} into a scratch
copy of an audited theorem --- must cause the gate to refuse; a gate that
cannot fail certifies nothing.

\section{praxis-graphlaw: Executable Law-State Evaluation}
\label{sec:praxis-graphlaw}

The \textsf{praxis-graphlaw} engine (version v26.7.5) is an executable law-state engine that evaluates the operational conformance of transition systems and process models. Rather than operating as a custom verification engine built from scratch, \textsf{praxis-graphlaw} is a derivative work extending the RDF/N3 reasoner \textsf{roxi} (originally developed by Bonte et al.). It is integrated into the generation and validation lifecycle of \textsf{ggen} through the \texttt{GraphEngine} abstraction boundary, specifically concrete implementations of the \texttt{GraphLawStore}.

The engine models law-states directly as Resource Description Framework (RDF) triple graphs, using Turtle (.ttl) format serialization. Operational logic and state transitions are defined using Notation3 (N3) rules and stratified Datalog, allowing for deterministic materialization, solution generation, and query evaluation via the reasoner's \texttt{materialize}, \texttt{prove}, and \texttt{solve} interfaces.

State conformance and safety properties are enforced via W3C SHACL (Shapes Constraint Language) and ShEx (Shape Expressions) constraint schemas, which operate as executable transaction-path admission gates. The engine exposes these gates via the \texttt{validate\_shacl} and \texttt{validate\_shex} functions. Safety violations and negative conditions are captured as N3 logical denials of the form \texttt{\{ body \} => false.}, where any deduction of falsity immediately triggers a transaction refusal.

At the integration layer with \textsf{ggen}, graphlaw's validation results are translated into a structured, typed refusal vocabulary ranging from \texttt{FM-LAW-001} through \texttt{FM-LAW-013}. For example, an N3 logical denial violation produces an \texttt{FM-LAW-011} refusal, while a SHACL shapes non-conformance results in an \texttt{FM-LAW-013} refusal. This structured failure reporting constitutes an independent, executable instance of the MathProofOps law-state pattern.

Fast evaluation is a critical design requirement for these gates, as they sit directly on the transaction-path execution loop of the process engine. High latency during validation would directly degrade transaction throughput. While the primary objective is low latency, we measure and record these performance characteristics empirically in Section~\ref{sec:rslab_benchmark} to verify that the overhead of validation remains within acceptable limits.

\section{rslab: Empirical Evidence Rail}
\label{sec:rslab}

The \textsf{rslab} framework acts as the empirical evidence rail for MathProofOps. It sits in parallel with the formal proof and correspondence rails, ensuring that measured performance and test execution data are receipted and ledgered deterministically. The governance of this rail is defined by the following doctrine:

\begin{quote}
\itshape
praxis executes. \\
rslab measures and receipts. \\
mfact admits formal standing. \\
paper renders admitted/receipted claims.

rslab is not a proof engine. \\
rslab is an empirical evidence rail.
\end{quote}

Within this rail, raw telemetry and outputs from execution are treated as unadmitted observations ($O$). Once these observations are processed by validation scripts, checked against JSON schemas, and bound to a cryptographic content receipt, they become schema-validated admitted observations ($O^*$). Downstream LaTeX paper fragments are then generated directly from these receipts to produce the final paper artifacts with standing ($A$).

Every experiment is declared in the central experiment manifest (\texttt{rslab/manifest.toml}) using the status ladder pattern. Telemetry outputs are validated against JSON schemas, such as \texttt{benchmark\_result.schema.json}. The resulting TOML receipts mirror the metadata structure of \texttt{verif-receipt.json}, containing the execution builder, toolchain specifications (e.g. pinned compiler versions), git commits, and BLAKE3 hashes of raw captures. Crucially, in accordance with the determinism requirements, these receipts contain no wall-clock timestamps in their body.

The fragment generation pipeline is strictly fail-closed: if the processed results file or the validating TOML receipt is missing or has a hash mismatch, the builder script exits non-zero with an \texttt{RSLAB\_EVIDENCE\_MISSING} refusal, halting the build rather than allowing blank or stale placeholders to be compiled.

Currently, \textsf{rslab} defines a readiness contract capturing throughput and latency metrics for core graphlaw operations. However, profiling data is not available. This limitation is a direct reflection of the development environment: no profiling or flamegraph tooling exists in the \textsf{praxis} workspace. This absence is explicitly declared in the receipt caveats rather than silently ignored.

We explicitly distinguish \textsf{rslab}'s role from formal verification. Where the correspondence rail (Section~\ref{sec:correspondence}) proves a bounded relationship between an implementation and admitted semantics, the \textsf{rslab} rail records what a system measurably does, without claiming that measurement constitutes proof of anything beyond itself.


\section{Use of Generative AI Tools}
\label{sec:ai-disclosure}

Claude Code and related language-model tooling were used to generate
candidate patches, proof sketches, ontology fragments, documentation
drafts, and release metadata during the run described in
Section~\ref{sec:run}. These outputs were treated throughout as untrusted
candidate artifacts, per the separation in Section~\ref{sec:mfact}: they
were admitted into \textsf{procint} only when accepted by the Lean kernel,
the sorry census, the axiom audit, schema validation against the
\texttt{procint:} ontology, and the negative-fixture and release-gate
checks of Section~\ref{sec:run}. No language-model output is part of the
trusted base (Table~\ref{tab:stages}); the trusted base is the Lean kernel,
Mathlib and CSLib at their pinned revisions, and the \texttt{lake}/\texttt{ggen}
toolchain. The author reviewed and accepts responsibility for all submitted
content, code, citations, and claims in this paper and in the accompanying
repository.

This separation is not precautionary in the abstract: general-purpose
language models have a documented Lean~3/Lean~4 syntax-confusion failure
mode, generating deprecated Lean~3 constructs despite explicit Lean~4
instructions, and GPT-4-class and substantially smaller models have been
reported at comparable, low accuracy on the MiniF2F-Valid
benchmark~\cite{langconfusion2025}. Approaches that fine-tune a
general-purpose model toward Lean~4 specifically~\cite{theoremllama2024} or
pair a language model with solver-guided refinement~\cite{apollo2025} trade
some of this untrusted-candidate distance for higher raw acceptance rates;
\textsf{mfact} does not take that trade in the run reported here, since every
candidate --- however produced --- passes through the same kernel-and-audit
gate regardless of its source's measured reliability.

\section{Implementation Correspondence with Aeneas}
\label{sec:correspondence}

Beyond \textsf{procint}'s own catalog, \textsf{mfact} opens a second
manufacturing rail: implementation correspondence. Where
Section~\ref{sec:procint} manufactures domain semantics from a declared
obligation graph, this rail manufactures obligations that relate a Lean
declaration to code extracted from an \emph{external}, unverified Rust
crate (\texttt{wasm4pm-core}) via Aeneas~\cite{ho2022aeneas} and its
Charon front end~\cite{protzenko2024charon}. The first instance of this
rail is deliberately scoped to one obligation --- token-replay counts
correspondence, ``D1'' --- not because the rail is narrow in principle,
but because a single bounded instance is what standing requires before the
rail generalizes. It follows the same status ladder as the rest of this
paper's standing (declared, extracted, stated, proven), computed by a
dedicated builder (\texttt{scripts/build\_verif.py}) from observed
evidence: a charon/aeneas pipeline receipt, a \texttt{lake build} of the
extracted module against \textsf{procint}, and an axiom-print check on the
resulting theorem. The
crate's \texttt{ReplayCounts}\allowbreak\texttt{\{produced, consumed,
missing, remaining\}} counts and its integer
\texttt{fitness\_num}/\texttt{fitness\_den} functions are extracted
verbatim; an abstraction function lifts the extracted (proof-free) struct
into \textsf{procint}'s \texttt{ReplayCounts} (which carries
\texttt{missing} $\le$ \texttt{consumed} and \texttt{remaining} $\le$
\texttt{produced} as proof fields) given those two inequalities as
hypotheses. The correspondence theorem does not assert that the crate's
integer ratio and \textsf{procint}'s exact-rational fitness are equal ---
they are structurally different formulas (a single integer fraction versus
an average of two ratios) --- only that each computes correctly from the
same underlying counts; forcing a false numeric equality between them was
rejected in favor of this honest, weaker, but PROVEN statement.
Table~\ref{tab:correspondence} is rendered by that same builder from its
own receipt; it is not written by hand.

\ifreleasebuild
  % AUTO-RENDERED by scripts/build_verif.py — do not edit by hand.
% Regenerate: python3 scripts/build_verif.py
\begin{table}[h!]
\centering
\begin{tabular}{lll}
\toprule
\textbf{Obligation} & \textbf{Status} & \textbf{Evidence} \\
\midrule
token\_replay\_counts\_corr & \textcolor{green}{PROVEN} & full ladder evidence present (extracted, stated, proven) \\
\bottomrule
\end{tabular}
\caption{Correspondence-factory obligation status, computed from observed charon/aeneas/lake evidence by \texttt{scripts/build\_verif.py} (status ladder: DECLARED $<$ EXTRACTED $<$ STATED $<$ PROVEN; never hand-set past DECLARED). See release/verif-receipt.json for the full evidence trail.}
\label{tab:correspondence}
\end{table}

\else
  \InputIfFileExists{correspondence_status}{}{%
  \begin{center}\itshape Draft build: correspondence\_status.tex missing.\end{center}}
\fi

\section{Evaluation}
\label{sec:eval}

\subsection{Formal Standing Evaluation}
\label{sec:formal_eval}

Every number in this section is rendered from the release manifest
(\texttt{release-manifest.json}) of the run; none is written by hand. The
generation pipeline fails closed: if the manifest is absent or stale, this
section does not render and the paper does not build for release.

\ifreleasebuild
  % GENERATED from release/release-manifest.json,
% release/gates.json, and `git log` on this repository.
% Extraction commands:
%   jq '.artifacts | length' release-manifest.json
%   jq '[.artifacts[] | select(.proven)] | length' release-manifest.json
%   jq '.statedNotProven' release-manifest.json
%   jq -r '.artifacts[] | select(.proven) | .axioms | sort | join(",")' release-manifest.json | sort | uniq -c
%   jq '.foldHash' release-manifest.json
%   find procint/ProcInt -name '*.lean' | xargs wc -l | tail -1
%   git log --format='%aI %s'
% Do not edit by hand. If this file is stale relative to release-manifest.json,
% regenerate before submission.

\begin{table}[h]
\centering
\small
\begin{tabular}{lr}
\toprule
Metric & Value \\
\midrule
Lines of Lean (\texttt{ProcInt/}) & 11{,}269 \\
Declarations recorded in the ontology & 397 \\
Modules & 53 \\
Theorems kernel-admitted and axiom-audited (\emph{proven}) & 197 \\
Statements formalized but not discharged (\emph{stated}) & 7 \\
Definitions, structures, and test oracles & 193 \\
\bottomrule
\end{tabular}
\caption{Corpus size and the proven/stated/definition split, as recorded
in the release manifest.}
\label{tab:corpus}
\end{table}

\begin{table}[h]
\centering
\small
\begin{tabular}{lr}
\toprule
Axiom footprint of proven theorems & Count \\
\midrule
No axioms & 37 \\
\texttt{propext} & 28 \\
\texttt{Quot.sound}, \texttt{propext} & 26 \\
\texttt{Classical.choice}, \texttt{Quot.sound}, \texttt{propext} & 106 \\
\midrule
Total proven & 197 \\
Unauthorized axioms found & 0 \\
Transitive \texttt{sorryAx} found & 0 \\
\bottomrule
\end{tabular}
\caption{Axiom-audit result over all 197 proven theorems. Every set is a
subset of the trusted \{\texttt{propext}, \texttt{Classical.choice},
\texttt{Quot.sound}\}; \texttt{no\_valid\_objection}
(Section~\ref{sec:mfact}) is the strongest case, depending on none of them.}
\label{tab:axioms}
\end{table}

\paragraph{The two stated (unproven) declarations.} Both are formal
statements, not theorems: \texttt{WfNet.sound\_iff\_short\-Circuit\_live\_bounded\_statement}
(the crown-jewel target, Section~\ref{sec:procint}) and
\texttt{BranchingProcess.isUnfoldingOf\_statement} (an unfolding-correctness
statement in the Petri-net seed lane). Per
\texttt{crownJewel\_status}, the crown-jewel target's release-time status
is \texttt{"stated"}: neither direction of the iff is discharged in this
release. Two supporting lemmas about \texttt{WfNet.Sound}
(\texttt{reaches\_final}, \texttt{enabled\_of\_transition}) \emph{are}
proven and audited (Table~\ref{tab:axioms}).

\paragraph{Fixtures.} \texttt{ProcInt.Fixtures.Positive} and
\texttt{ProcInt.Fixtures.Negative} both build; the negative fixtures are
\texttt{\#guard\_msgs (error)} blocks whose docstring is the exact error
Lean produces for a malformed \texttt{ReplayCounts} (missing token count
exceeding consumed) and a non-monotone trace, so the corpus fails to build
if either rejection ever silently stops happening.

\paragraph{Certification.} \texttt{mfact certify release-manifest.json
gates.json} exits 0 with \texttt{sorryFree}, \texttt{axiomsClean},
\texttt{fixturesPass}, and \texttt{evidenceComplete} all true. The
genesis-folded release hash (BLAKE3, folded over all 397 artifact hashes
in name order, seed \texttt{"mfact-v26.7.7-genesis"}) is
\texttt{942facf3...f2d434} (truncated; full value in
\texttt{release-manifest.json}). Both negative controls behaved as
required: flipping \texttt{sorryFree} to false in a copy of the gates
file causes \texttt{mfact certify} to exit 1 with an explicit gate-failure
message; a syntactically corrupted manifest causes it to exit 2 with a
typed refusal, not a crash.

\paragraph{Wall-clock.} This section reports what the repository's git
history actually timestamps: from the paper-scaffold commit to the
certified-release commit, six commits span
\texttt{2026-07-06T20:41:32-07:00} to \texttt{2026-07-06T21:09:48-07:00}
(28 minutes), covering discovery and repair of three cross-family
import/duplicate-declaration bugs, addition of the fixture pair, the
axiom-audit reconciliation pass over 145 theorems, and certification. The
ontology authoring and initial corpus rendering that preceded the first
commit in this span are not independently timestamped in this
repository's history; we report the receipted span only and do not
extrapolate a total-run duration from it.

\else
  \InputIfFileExists{evaluation}{}{%
  \begin{center}\fbox{\itshape Draft build: evaluation.tex missing.}\end{center}}
\fi

\paragraph{Final status.} The table below is generated from the
release manifest and standing records; no cell is hand-authored.
\ifreleasebuild
  \input{final_status}
\else
  \InputIfFileExists{final_status}{}{%
  \begin{center}\fbox{\itshape Draft build: final\_status.tex missing.}\end{center}}
\fi

\subsection{Standing Quadrature}
\label{sec:quadrature}

The release is not evaluated only by whether Lean builds. It is evaluated
by whether the declaration catalog, admitted Lean corpus, process
evidence, release manifest, and paper claims form a \emph{closed standing
graph}. The quadrature check rejects orphan declarations, unsupported
claims, untraced artifacts, and manifest--paper divergence. Closure is
itself kernel-checked: a generated witness module
(\texttt{ProcInt.Release.Quadrature}) carries the name-sorted proven
surfaces of the catalog, the manifest, and the axiom audit as literal
lists, and proves their pairwise equality by \texttt{rfl} --- any orphan
on any side makes the lists differ and the kernel refuses the witness.
The same module encodes the manufacturing run as a \textsf{procint}
process trace and proves, with the corpus's own \texttt{DeclareConstraint}
semantics, that the run conforms to its declared process model:
\textsf{procint} models the process by which \textsf{mfact} manufactured
\textsf{procint}. Table~\ref{tab:quadrature} is generated from
\texttt{quadrature.json}; three negative controls (a corrupted evaluation
number, a claim with no evidence, a catalog declaration with no rendered
symbol) each cause a typed refusal.

\ifreleasebuild
  % RENDERED by ggen sync from the quadrature graph (quadrature-pack).
% Do not edit: regenerate via `just standing-quadrature`.
\begin{table}[h]
\centering
\small
\begin{tabular}{lrl}
\toprule
Surface pair & Count & Result \\
\midrule
TTL to Lean & 197 & PASS \\
Lean to Audit & 197 & PASS \\
Audit to Manifest & 197 & PASS \\
Manifest to Paper & 6 & PASS \\
Artifact to Process Event & 254 & PASS \\
Claim to Evidence & 5 & PASS \\
\bottomrule
\end{tabular}
\caption{Standing Quadrature: PASS over
197 proven declarations, 5 traced paper
claims, and 6 evaluation numbers. The closure is
kernel-checked by the rendered witness
\texttt{ProcInt.Release.Quadrature}.}
\label{tab:quadrature}
\end{table}

\else
  \InputIfFileExists{quadrature}{}{%
  \begin{center}\itshape Draft build: quadrature.tex missing.\end{center}}
\fi

\subsection{praxis/rslab Benchmark Evidence}
\label{sec:rslab_benchmark}

\ifreleasebuild
  \input{../rslab/paper_fragments/rslab_praxis_graphlaw_summary}
\else
  \InputIfFileExists{../rslab/paper_fragments/rslab_praxis_graphlaw_summary}{}{%
  \begin{center}\itshape Draft build: rslab\_praxis\_graphlaw\_summary.tex missing.\end{center}}
\fi

\ifreleasebuild
  \input{../rslab/paper_fragments/rslab_praxis_graphlaw_benchmarks}
\else
  \InputIfFileExists{../rslab/paper_fragments/rslab_praxis_graphlaw_benchmarks}{}{%
  \begin{center}\itshape Draft build: rslab\_praxis\_graphlaw\_benchmarks.tex missing.\end{center}}
\fi

\ifreleasebuild
  \input{../rslab/paper_fragments/rslab_praxis_graphlaw_profiles}
\else
  \InputIfFileExists{../rslab/paper_fragments/rslab_praxis_graphlaw_profiles}{}{%
  \begin{center}\itshape Draft build: rslab\_praxis\_graphlaw\_profiles.tex missing.\end{center}}
\fi

\ifreleasebuild
  \input{../rslab/paper_fragments/rslab_readiness}
\else
  \InputIfFileExists{../rslab/paper_fragments/rslab_readiness}{}{%
  \begin{center}\itshape Draft build: rslab\_readiness.tex missing.\end{center}}
\fi

\section{Falsifier and Valid Objection Surface}
\label{sec:falsifier}

A valid objection to this release must exhibit at least one of:
(1) a declared admitted theorem contains \texttt{sorry};
(2) a theorem depends on unauthorized axioms;
(3) a manifest hash does not match the artifact;
(4) a negative fixture passes when it should fail;
(5) a generated evaluation number does not match the manifest;
(6) a stated declaration is described as proven;
(7) the catalog-to-Lean rendering changes the intended theorem statement;
(8) the release cannot be replayed under pinned dependencies.
These are \emph{closed failure classes}, not open attack surfaces: each
is engineered into a gate, a typed refusal, or a non-admitted state, and
the release gate is designed so that any occurrence of such a case
prevents admission or produces a typed refusal. After certification, no
constructor of the valid-objection surface is inhabited inside the
declared boundary; the negative controls demonstrate each class refusing
on a poisoned copy. A local objection to one proof, citation, or template
is not automatically a refutation of mathematical manufacturing. It
becomes a refutation only if it crosses the admission boundary and
invalidates a claimed receipt --- and the enumeration above is exactly
the set of ways to do so.

\section{Limitations and Standing}
\label{sec:limitations}

\paragraph{Proven versus stated.}
The library's standing is exactly the manifest's proven/stated split; prose
in this paper adds nothing to it. \emph{Stated} declarations are formal
statements whose proof obligations are open --- they compile, they are
citable as definitions of the intended property, and they carry no more
authority than that.

\paragraph{Crown-jewel scope.}
The short-circuit soundness characterization is reported per direction. If
only one direction is admitted at release time, the claim is that direction
and nothing more; the converse remains a stated obligation in the manifest.
\ifreleasebuild
  % GENERATED FILE — DO NOT EDIT BY HAND (source: quadrature graph via ProcInt.crownJewel_status)
The release-time classification of the short-circuit soundness
characterization is \textbf{\CrownJewelStatus} (ground truth:
\texttt{ProcInt.crownJewel\_status}); this fragment is generated from the
catalog and reports exactly that ground-truth classification, whichever
value it holds — never a hand-asserted status.

\else
  \InputIfFileExists{crown_jewel_status}{}{}
\fi

\paragraph{What admission does not check.}
Kernel admission plus the axiom audit establishes that admitted terms prove
the stated theorems under the pinned axioms. It does not establish that
the catalog-curated statements faithfully render the informal canon; that
correspondence rests on citation-annotated statement curation against the
primary sources and is exactly the kind of residual risk the semantic
audits of~\cite{meek2026formalizing} target. Fidelity of statements is a
review obligation, not a theorem.

\paragraph{Scope of the novelty claim.}
``First mechanization'' claims are scoped to a search of the Lean
(Mathlib/CSLib), Coq, and Isabelle (AFP) ecosystems at run time; we claim a
negative search result, not a proof of absence.

\paragraph{Standing at submission.}
The standing of the system is intentionally not inferred from this paper.
At release time, the crown theorem rail, countermodel rail, correspondence
rail, product/DX rail, and paper rail each carry independent manifest
status. The paper is a publication surface, not a source of truth; if
paper prose and manifest status ever diverge, the manifest and receipts
govern, not the prose.

\paragraph{Exclusions.}
This paper does \emph{not} claim any of the following.

\begin{center}
\small
\begin{tabular}{ll}
\toprule
Exclusion & Reason \\
\midrule
LLMs prove mathematics autonomously & LLMs are untrusted producers \\
\textsf{procint} formalizes all of process intelligence & v1 boundary is manifest-scoped \\
Kernel admission proves informal intent & statement fidelity is a review obligation \\
Rice's theorem is violated & arbitrary semantic correctness is not decided \\
External opinion enters the standing calculus & only receipted artifacts are parsed \\
Criticism outside the boundary is refuted & the empty-objection-surface claim is scoped to the declared boundary \\
\textsf{ggen} output has standing by generation & only admitted output has standing \\
A stated theorem is proven & stated means the obligation remains open \\
\bottomrule
\end{tabular}
\end{center}

\section{Availability}
\label{sec:availability}

The \textsf{mfact} and \textsf{procint} packages, the ontology, the
templates, the release manifest, and the source of this paper (including
the generated evaluation) will be made available in the \texttt{mfact}
repository at the certified release tag accompanying this submission,
pinned to Lean~4 v4.31.0 with Mathlib and CSLib at fixed revisions. The
negative-control script that demonstrates the gate refusing a
\texttt{sorry}-injected artifact ships with the release.

\paragraph{Reproducibility appendix.}
\begin{itemize}
\item \textbf{Lean:} \texttt{leanprover/lean4:v4.31.0}.
\item \textbf{Mathlib:} \texttt{leanprover-community/mathlib4} at the commit
  tagged \texttt{v4.31.0}.
\item \textbf{CSLib:} \texttt{leanprover/cslib}, pinned to its last commit
  on the \texttt{v4.31.0} toolchain.
\item \textbf{Build:} \texttt{lake build} in \texttt{mfact/} and
  \texttt{procint/}.
\item \textbf{Rendering:} \texttt{ggen sync run} against
  \texttt{packs/lean-math-pack} and the merged \texttt{procint:} catalog.
  This renders the candidate Lean corpus (modules, index, axiom audit,
  registry tables) from the TTL catalog; the rendered corpus is admitted
  only by the subsequent \texttt{lake build}.
\item \textbf{Certification:} \texttt{lake exe mfact certify
  release-manifest.json gates.json}.
\item \textbf{Artifact and manifest hashes:} recorded per-declaration
  (BLAKE3) and as a genesis-folded release hash in
  \texttt{release-manifest.json}; reproduced in the generated evaluation
  (Section~\ref{sec:eval}).
\item \textbf{Expected output:} \texttt{CERTIFIED\_RELEASE=PASS} with
  \texttt{SORRY\_COUNT=0} and \texttt{AXIOM\_AUDIT=PASS}; a negative control
  with an injected \texttt{sorry} must instead exit refused.
\end{itemize}

\ifreleasebuild
  % GENERATED FILE — DO NOT EDIT BY HAND (source: quadrature graph; see release_macros.tex)
\paragraph{Replay.}
The release is tag \texttt{\ReleaseTag} (release hash
\texttt{\footnotesize\ReleaseHash}). A clean replay regenerates the corpus
and re-certifies it:
\begin{center}
\small
\begin{tabular}{ll}
\toprule
Step & Command \\
\midrule
Render corpus from catalog & \texttt{just render} \\
Build + audits (both packages) & \texttt{just build \&\& just audit} \\
Regenerate manifest + gates & \texttt{just manifest} \\
Certify (exit 0 required) & \texttt{just certify} \\
Standing Quadrature & \texttt{just standing-quadrature} \\
Quadrature negative controls & \texttt{just quadrature-negative-controls} \\
arXiv package & \texttt{just arxiv-package} \\
\bottomrule
\end{tabular}
\end{center}
The expected result is not ``trust us''; it is a receipt: certification
exits 0, the quadrature witness kernel-admits, and the negative controls
refuse.

\else
  \InputIfFileExists{availability}{}{}
\fi

\ifreleasebuild
  % RENDERED by `ggen sync` from the post-release graph (post-release-pack) — DO NOT EDIT BY HAND
\paragraph{Publication packet.}
Publication of v26.7.7 is manufactured as its own identity, distinct
from the certified core release: a packet hash folded over the packet's
evidence inputs (recorded in \texttt{release/final\_status.json}; this
paper is itself one of those inputs and therefore cannot embed the hash
without self-reference), with each actuation surface gated on evaluated
requirements. The arXiv packet is
\texttt{PENDING} (first pass (--plan): fragments not yet rendered;
13 declared files, declared $=$ packed enforced).
Packet statuses:
\texttt{arxiv\_upload}$=$\texttt{BLOCKED}, \texttt{github\_push}$=$\texttt{BLOCKED}, \texttt{github\_release}$=$\texttt{ALIVE}.
No packet self-actuates: publication is recorded as
\texttt{PENDING\_EXTERNAL\_ACTUATION} in every case, and the
kernel-admitted witness \texttt{ProcInt.Release.PostRelease} refuses any
render in which a packet claims otherwise.


  % RENDERED by `ggen sync` from the post-release graph (post-release-pack) — DO NOT EDIT BY HAND
\paragraph{Independent replay.}
The replay lane is upgraded only by its own gate's report, never asserted:
its current standing is \texttt{REPLAY\_PASS}. The
gate clones the repository at tag \texttt{v26.7.7-procint-certified}
into a scratch tree and re-runs the manufacturing rail there:
\begin{enumerate}
  \item \texttt{git clone \textless stand-in\textgreater /mfact-dryrun.git at tag v26.7.7-procint-certified}
  \item \texttt{just render \&\& just build (pinned toolchain, lake exe cache get)}
  \item \texttt{just audit \&\& just test}
  \item \texttt{just certify (exit 0 required)}
  \item \texttt{just regen-check (ARTIFACT\_DRIFT\_REFUSED on any divergence)}
\end{enumerate}
A replay that diverges on any ledgered artifact refuses with
\texttt{ARTIFACT\_DRIFT\_REFUSED}; a replay that has not been run reports
\texttt{REPLAY\_NOT\_RUN} rather than borrowing the core release's standing.


\else
  \InputIfFileExists{publication_status}{}{}
  \InputIfFileExists{replay_status}{}{}
\fi

\section{Conclusion}

\ifreleasebuild
  % GENERATED FILE — DO NOT EDIT BY HAND (source: quadrature graph; see release_macros.tex)
We have presented \textsf{mfact} as a Lean~4 framework for receipted
mathematical manufacturing and \textsf{procint} as its first manufactured
process-intelligence domain. The certified release \texttt{\ReleaseTag}
contains \KernelTheorems{} kernel-proven and axiom-audited theorems out of
\TotalDecls{} declarations, with \SorryCount{} sorries, positive and
negative fixture coverage, and a BLAKE3-folded release hash. The central
claim is not that generation is trustworthy; it is the opposite:
generation is untrusted, and therefore mathematical standing must be
manufactured through explicit admission, refusal, audit, manifest, and
replay.

This changes the unit of formalization. A theorem is no longer treated
merely as a local proof term. In a manufactured release, a theorem is an
admitted artifact with provenance, hash, axiom footprint, fixture status,
manifest position, and replay path. A failed theorem is not erased or
softened; it is refused or downgraded to stated status. A release is not
declared successful by prose; it is certified by a receipt-bearing
boundary. The release also closes Standing Quadrature
(\QuadratureStatus): the declaration catalog, manifest, and audit
surfaces, the claim--evidence ledger, and the process-trace conformance
theorems are connected by generated witnesses admitted by Lean.

The crown-jewel workflow-net theorem is classified as
\emph{\CrownJewelStatus}: the short-circuit
characterization of soundness by liveness and boundedness
(van der Aalst 1997, Lemma~8 / Theorem~11) is a kernel-admitted,
axiom-audited theorem in this release, assembled from independently
proven directions rather than asserted. This is
not a release failure; it is the release boundary doing its job either way.

The broader result is a bounded response to the problem of generated
formal content. Rice-style undecidability is not denied; arbitrary
semantic judgment is replaced by declared admission inside a replayable
boundary. Contact is not consequence; candidate generation is not proof;
compilation is not full standing; and prose is not receipt. A
mathematical artifact stands only when \(R_B \vdash A = \mu(O^*_B)\).
The release therefore does not ask for belief; it eliminates belief as a
standing mechanism.

\else
  \InputIfFileExists{conclusion}{}{%
  \begin{center}\fbox{\itshape Draft build: conclusion.tex missing.}\end{center}}
\fi

\bibliographystyle{plain}
\bibliography{refs}

\end{document}
